% =========================================================
% Proof Stub: Riemann Rearrangement Theorem
% Source: volume-iii/analysis/sequences/notes/series/notes-series-rearrangements.tex
% =========================================================

% *** HARDER PROOF — attempt after foundational results settled ***

\clearpage
\phantomsection
\hypertarget{proof-RA-ABB-T-riemann-rearrangement}{}

\subsubsection[Riemann Rearrangement Theorem]{Proof --- Riemann Rearrangement Theorem}

\bigskip

\noindent
\textbf{Source.} See \texttt{notes-series-rearrangements.tex}.

\vspace{0.75em}

\noindent
\textbf{Goal.} If $\sum a_n$ converges conditionally, then for any $L \in \mathbb{R}$ (or $L = \pm\infty$), there is a rearrangement converging to $L$.

\vspace{0.75em}

\noindent
\textbf{Logical form.}
\[
  \sum a_n \text{ cond. conv.} \Rightarrow \forall L \in \overline{\mathbb{R}},\; \exists \sigma,\; \sum a_{\sigma(n)} = L.
\]

\vspace{0.75em}

\noindent
\textbf{Proof strategy.} The positive and negative parts $p_n$ and $q_n$ each diverge (since conditional convergence $+$ absolute divergence forces this). Greedily pick enough positive terms to exceed $L$, then enough negative terms to go below, alternating; the partial sums oscillate around $L$ with diminishing swings since $a_n \to 0$.

\vspace{0.75em}

\noindent
\textbf{Proof.}
\begin{proof}
\Claim If $\sum a_n$ converges conditionally, then for any $L \in \mathbb{R}$ (or $L = \pm\infty$), there is a rearrangement converging to $L$.

\Given 

\Goal 

\Strategy The positive and negative parts $p_n$ and $q_n$ each diverge (since conditional convergence $+$ absolute divergence forces this). Greedily pick enough positive terms to exceed $L$, then enough negative terms to go below, alternating; the partial sums oscillate around $L$ with diminishing swings since $a_n \to 0$.

\medskip

% --- let p_n = max(a_n,0), q_n = max(-a_n,0); both ∑p_n and ∑q_n diverge ---
% --- greedy construction: pick positives until partial sum exceeds L; then pick negatives until it goes below L; repeat ---
% --- each overshoot is bounded by the next term |a_n|; since a_n → 0, the oscillations → 0 ---
% --- partial sums converge to L ---

\end{proof}

\begin{remark*}[Proof shape]
Greedy construction exploiting the divergence of both positive and negative parts; the oscillation shrinks because $a_n \to 0$.
\end{remark*}

\begin{remark*}[Dependencies]
\begin{itemize}
  \item $a_n \to 0$ (necessary for any convergent series).
  \item Both $\sum p_n = +\infty$ and $\sum q_n = +\infty$ (conditional convergence).
  \item Greedy selection algorithm.
\end{itemize}
\end{remark*}
